% !PRINTED 706
\weberpage{706}
Um nun diesen Ausdruck für die Function $\varPhi$ weiter umzuformen
und zur Berechnung vorzubereiten, hat man die Sätze
über die \gesperrt{Abel}'schen Gruppen und Gruppencharaktere zu benutzen,
die wir in den Paragraphen 14 bis 16 dieses Bandes
kennen gelernt haben.

Die Gesammtheit der durch $q$ nicht theilbaren Zahlen $n$
bildet, nach dem Modul $m$ genommen, eine \gesperrt{Abel}'sche Gruppe $\mathfrak{N}$
vom Grade $\mu$, deren Charaktere $\chi(n)$ im \S.\,16 bestimmt sind. Es
war dabei ein kleiner Unterschied, je nachdem $q$ ungerade oder
$q=2$ ist.

1) Ist $q$ ungerade, so giebt es eine primitive Wurzel $c$ von $m$,
und für jede Zahl $n$ giebt es einen Index $\gamma$, so dass
\weq{5}{n \equiv c^{\gamma}\ (\mathrm{mod}\ m).}
Ist dann $\Theta$ eine $\mu^{\text{te}}$ Einheitswurzel, so ist
\weq{6}{\chi(n) = \Theta^{\gamma},}
und die sämmtlichen $\mu$ Charaktere erhält man, wenn man für $\Theta$
die verschiedenen $\mu^{\text{ten}}$ Einheitswurzeln setzt.

Gehört $n$ zum Exponenten $f$, so ist $e$ der grösste gemeinschaftliche
Theiler von $\gamma$ und $\mu$, und $\chi(n)$ ist $f^{\text{te}}$ Einheitswurzel.

2) Für $q=2$, $m>4$ (den Fall $m=4$ berücksichtigen wir
hier nicht) hat jede ungerade Zahl $n$ zwei Indices $\alpha, \beta$, die nach
den Moduln $2, \frac{1}{2}\mu$ bestimmt sind, für die
\weq{7}{n \equiv (-1)^{\alpha}\,5^{\beta}\ (\mathrm{mod}\ m)}
ist. Als Charaktere erhält man, wenn $\varepsilon = \pm 1$ und $\Theta$ gleich
einer $\frac{1}{2}\mu^{\text{ten}}$ Einheitswurzel gesetzt wird:
\weq{8}{\chi(n) = \varepsilon^{\alpha}\Theta^{\beta};}
$f$ ist die kleinste positive Lösung der beiden Congruenzen:
\[
\alpha f \equiv 0\ (\mathrm{mod}\ 2), \quad \beta f \equiv 0\ (\mathrm{mod}\ \tfrac{1}{2}\mu),
\]
und $\chi(n)$ ist auch hier $f^{\text{te}}$ Einheitswurzel.

Für beide Fälle gilt aber die Bemerkung:

3) Gehört $n$ zum Exponenten $f$, so erhält man, wenn man $\chi$
die gesammte Gruppe der $\mu$ Charaktere durchlaufen lässt, aus
$\chi(n)$ jede $f^{\text{te}}$ Einheitswurzel $e$ mal.

Denn da die $\chi$ eine \gesperrt{Abel}'sche Gruppe bilden, so erhält
man zunächst jede $f^{\text{te}}$ Einheitswurzel, die überhaupt vorkommt,
gleich oft, nämlich so oft als den Werth 1, weil, wenn $\chi_1(n)=\chi_2(n)$
ist, $\chi_1\chi_2^{-1}(n) = \chi_0(n) = 1$ ist. Andererseits erhält man aber
%CONT>>
